\section*{الفصل \ref{ch:g11:absorbance} --- اللون والامتصاصية}

\begin{solution}{exo:g11:absorbance:1}
إذا امتص البرتقالي بدا أزرق، وهو اللون المقابل للبرتقالي على الدائرة.
وإذا امتص البنفسجي بدا أصفر.
\end{solution}

\begin{solution}{exo:g11:absorbance:2}
\omterm{def:g2:mixing-and-dissolving:dissolve}{المحلول} الأخضر يمتص أساسًا الضوء الأحمر، وهو اللون المقابل للأخضر على
الدائرة: من نحو \qty{620}{nm} إلى \qty{750}{nm}.
\end{solution}

\begin{solution}{exo:g11:absorbance:3}
الصبغة الزرقاء: $\lambda_{\max} = \qty{629}{nm}$، $A = 0.82$. الصبغة الصفراء:
$\lambda_{\max} = \qty{427}{nm}$، $A = 0.795$، أي نحو 0.80.
\end{solution}

\begin{solution}{exo:g11:absorbance:4}
$A$ تتناسب مع التركيز: فثلاثة أضعافه تعطي
$A = 0.90$؛ ونصفه يعطي $A = 0.15$.
\end{solution}

\begin{solution}{exo:g11:absorbance:5}
الشاهد خلية مملوءة \omterm{def:g6:solutions-and-solubility:solute}{بالمذيب} النقي. وضبط $A = 0$ به يُلغي ما تمتصه الخلية \omterm{def:g6:solutions-and-solubility:solute}{والمذيب}
والجهاز، فتكون القراءات التالية راجعة إلى \omterm{def:g6:solutions-and-solubility:solute}{المذاب}
وحده.
\end{solution}

\begin{solution}{exo:g11:absorbance:6}
$C = 0.41 / 0.164 = \qty{2.5}{mg/L}$.
\end{solution}

\begin{solution}{exo:g11:absorbance:7}
عند \qty{629}{nm} تكون \omterm{def:g11:absorbance:absorbance}{الامتصاصية} أكبر ما يمكن، فيكون القياس أشد حساسية (فالتراكيز
الصغيرة لا تزال تعطي امتصاصيات قابلة للقراءة)؛ وعند قمة الحزمة يكون المنحنى
مسطحًا، فلا يكاد خطأ صغير في طول الموجة يغيّر القراءة. أما عند \qty{550}{nm}
\omterm{def:g11:absorbance:absorbance}{فالامتصاصية} صغيرة وتتغير بسرعة مع طول الموجة.
\end{solution}

\begin{solution}{exo:g11:absorbance:8}
المخفف: $C = 0.48 / 0.164 = \qty{2.9}{mg/L}$؛ والمشروب:
$5 \times 2.9 = \qty{15}{mg/L}$ (وبدقة أكبر $14.6$).
\end{solution}

\begin{solution}{exo:g11:absorbance:9}
$a = A / (\ell\, C_m) = 0.795 / (1.00 \times 0.0150) = \qty{53.0}{L/(g.cm)}$;
$\varepsilon = a \times M = 53.0 \times 534.4 = \qty{2.83e4}{L/(mol.cm)}$.
\end{solution}

\begin{solution}{exo:g11:absorbance:10}
فوق نحو \qty{520}{nm}. وعند \qty{629}{nm} لا تمتص الصبغة الصفراء أبدًا،
\omterm{def:g11:absorbance:absorbance}{فالامتصاصية} هناك كلها راجعة إلى الصبغة الزرقاء.
\end{solution}

\begin{solution}{exo:g11:absorbance:11}
تتضاعف: فالقيمة $A = \varepsilon \ell c$ تتناسب مع $\ell$؛ والضوء يعبر ضعف
\omterm{def:g2:mixing-and-dissolving:dissolve}{المحلول}، أي ضعف \omterm{def:g7:atoms-and-molecules:molecule}{الجزيئات} الماصّة.
\end{solution}

\begin{solution}{exo:g11:absorbance:12}
$A/C$: $0.164$، $0.162$، $0.115$، $0.0675$. والنسبة ثابتة للأول والثاني فقط:
فبعد نحو \qty{10}{mg/L} (امتصاصيات فوق نحو 1.6) لا يعود القانون صحيحًا.
فلا يمكن استعمال إلا المحلولين المعياريين 5 و\qty{10}{mg/L} (وما دونهما)؛
والمجهول الذي يقرأ 2.5 يجب أن يُخفَّف، مثلًا خمس مرات، ويُقاس من
جديد.
\end{solution}

\begin{solution}{exo:g11:absorbance:13}
الصبغة الزرقاء، من \qty{629}{nm} وحدها: $C = 0.574 / 0.164 = \qty{3.5}{mg/L}$.
والصبغة الصفراء، من \qty{427}{nm}: $C = 0.53 / 0.0530 = \qty{10}{mg/L}$.
\end{solution}

\begin{solution}{exo:g11:absorbance:14}
تجد $0.368 / 0.164 = \qty{2.24}{mg/L}$. \omterm{def:g11:absorbance:absorbance}{والامتصاصية} الحقيقية
$0.368 - 0.040 = 0.328$، فالتركيز الحقيقي
$0.328 / 0.164 = \qty{2.00}{mg/L}$. فهي تبالغ بنسبة \qty{12}{\percent}.
\end{solution}

\begin{solution}{exo:g11:absorbance:15}
$10 \times 60 = \qty{600}{mg}$ في اليوم، في $600 / 20 = \qty{30}{L}$ من
المشروب. ولا أحد يشرب ثلاثين لترًا في اليوم: فالمشروب وحده ليس خطرًا واقعيًا،
وإن كانت الصبغة قد تأتي أيضًا من أغذية أخرى.
\end{solution}

\begin{solution}{pb:g11:absorbance:1}
\textbf{1.} البرتقالي الأحمر، وهو اللون المقابل للأزرق على الدائرة.

\textbf{2.} $\lambda_{\max} = \qty{629}{nm}$.

\textbf{3.} عند \qty{450}{nm} يكون الضوء أزرق: فالصبغة تمرره، وهذا بالضبط
سبب زرقة المشروب.

\textbf{4.} عند \qty{629}{nm}.

\textbf{5.} $\qty{50.0}{mg} / \qty{0.5000}{L} = \qty{100}{mg/L}$.

\textbf{6.} \omterm{def:g10:concentration-and-dilution:dilution}{معامل التخفيف} $100 / 3.0 = 33.3$، إذن
$V_0 = 100.0 / 33.3 = \qty{3.00}{mL}$: ماصّة مدرّجة (أو
سحاحة) وحوجلة معيارية من \qty{100.0}{mL}.

\textbf{7.} $A/C$ = 0.168، 0.163، 0.165، 0.163، 0.165: كلها قريبة من
0.164، إذن $A \approx 0.164\, C$، أي مستقيم يمر بالمبدأ.

\textbf{8.} \omterm{def:g2:mixing-and-dissolving:dissolve}{المحلول} الخالي من الصبغة ($C = 0$) هو الشاهد، المضبوط على
$A = 0$.

\textbf{9.} $a = \qty{0.164}{L/mg}$ لكل سنتيمتر، أي
\qty{164}{L/(g.cm)}؛ و$\varepsilon = 164 \times 792.9 =
\qty{1.30e5}{L/(mol.cm)}$.

\textbf{10.} $F = 50.0 / 25.0 = 2$.

\textbf{11.} $C = 0.590 / 0.164 = \qty{3.60}{mg/L}$.

\textbf{12.} $2 \times 3.60 = \qty{7.20}{mg/L}$.

\textbf{13.} نحو $2 \times 0.590 = 1.18$، فوق أعلى \omterm{def:g10:concentration-and-dilution:standard}{محلول معياري}
(0.823): فتقع القراءة خارج المجال المعاير، حيث قد لا يصح القانون. \omterm{def:g10:concentration-and-dilution:dilution}{والتخفيف}
يعيدها إلى ما بين المحاليل المعيارية.

\textbf{14.} $\qty{7.20}{mg/L} \times \qty{0.500}{L} = \qty{3.60}{mg}$.

\textbf{15.} $n = \num{3.60e-3} / 792.9 = \qty{4.54e-6}{mol}$.

\textbf{16.} $6 \times 30 = \qty{180}{mg}$ في اليوم.

\textbf{17.} $180 / 3.60 = 50$ قارورة.

\textbf{18.} $50 \times 0.500 = \qty{25}{L}$ في اليوم: يستحيل شربها. فصبغة هذا
المشروب وحدها لا يمكن أن تقرّب طفلًا من الحد، وإن كانت صبغات أغذية عدة
تتراكم.

\textbf{19.} $6 \times 70 = \qty{420}{mg}$، أي $420 / 3.60 \approx
117$ قارورة.

\textbf{20.} نحو 50 قارورة من \qty{500}{mL}، أي \qty{25}{L} من المشروب، في
يوم واحد.
\end{solution}
